\chapter{First-Order Multi-Pool (FOMP)}
\label{ch:fomp}

\section{Concept}

The First-Order Multi-Pool model simulates organic matter decomposition using two pools with different decay rates:

\begin{itemize}
  \item \textbf{Labile pool} --- Easily degradable fraction, fast decay ($k_{\text{labile}} \approx 0.1$--$1.0$/year)
  \item \textbf{Recalcitrant pool} --- Resistant fraction, slow decay ($k_{\text{recalcitrant}} \approx 0.01$--$0.05$/year)
\end{itemize}

Each pool follows first-order kinetics:

\begin{equation}
  \frac{dS_i}{dt} = I_i(t) - k_i \cdot S_i(t)
\end{equation}

where $S_i$ is the stock in pool $i$, $I_i$ is the inflow to pool $i$, and $k_i$ is the decay rate constant.

The analytical solution for each time step is:

\begin{equation}
  \text{Decay}_i(t) = S_i(t) \cdot \left(1 - e^{-k_i}\right)
\end{equation}

\section{Configuration (Sheet 3\_2)}

FOMP parameters are defined in the \texttt{3\_2\_Definition\_FOMP} sheet:

\begin{table}[h]
\centering
\begin{tabular}{llp{5.5cm}}
\toprule
\textbf{Parameter} & \textbf{Example} & \textbf{Description} \\
\midrule
\texttt{Process\_ID} & 4 & Process to apply FOMP to \\
\texttt{f\_labile} & 0.7 & Fraction of inflow to labile pool \\
\texttt{k\_labile} & 0.5 & Labile decay rate (1/year) \\
\texttt{k\_recalcitrant} & 0.025 & Recalcitrant decay rate (1/year) \\
\texttt{outflow\_id} & F\_04\_00\_CO2 & Carbon outflow (mineralization) \\
\texttt{outflow\_id\_2} & F\_04\_00\_env & Environmental outflow (non-carbon) \\
\bottomrule
\end{tabular}
\caption{FOMP parameter configuration.}
\end{table}

\begin{notebox}
The recalcitrant fraction is automatically calculated as $f_{\text{recalcitrant}} = 1 - f_{\text{labile}}$.
\end{notebox}

\section{Element Handling}

FOMP processes handle elements differently based on their nature:

\begin{description}
  \item[Carbon outflow] Contains the mineralized carbon mass. Elements: material~$=$~DM~$=$~CC~$=$~carbon mass.
  \item[Environmental outflow] Contains the non-carbon dry matter and bypassed water. Elements calculated as:
  \begin{itemize}
    \item WC = water content of inflow (bypasses decay)
    \item DM = non-carbon dry matter decay
    \item CC = 0 (all carbon goes to carbon outflow)
  \end{itemize}
\end{description}

\section{Use Cases}

FOMP is appropriate for modeling:
\begin{itemize}
  \item Composting processes (labile fraction decays in weeks, recalcitrant in years)
  \item Soil organic matter turnover
  \item Anaerobic digestion (with adjusted rate constants)
  \item Landfill degradation (very slow recalcitrant decay)
\end{itemize}

\begin{warnbox}
FOMP uses a time-averaged CC/DM ratio as a constant for element calculations. This approximation can cause a small mass balance error ($\sim$2\%) for the CC element. This is a known modeling limitation, not a bug.
\end{warnbox}
